\subsection{Tame weights on the actual quotient and the next exact height-two index}

\begin{lemma}[An equivariant coefficient bound for the actual target]
\label{lem:B-RD64-weighted-kernel}
\status{Candidate; complete coefficient and Cartier-module proof written.}
Let $p\ge5$, $N=p^2-1$, and let $L/K$ be a totally ramified
cyclic tame extension of degree $N$ of complete mixed-characteristic
$(0,p)$ discrete valuation fields with algebraically closed residue
field $k$. Put $R=\mathcal O_L$, normalize $v_L(\pi)=1$, and
write $a=v_L(p)$. Choose a generator $\tau$ with
$\tau\pi=\zeta\pi$, where $\zeta\in K$ has order $N$.
Let $G/R$ be a finite-height formal group with a semilinear
$\tau$-descent action. Assume $G_k$ is a product of one-dimensional
height-two formal groups, and its cotangent eigenvalues are
$\bar\zeta^{c_1},\ldots,\bar\zeta^{c_d}$ with
\[
                       c_i\equiv c_j\pmod{p-1}
                       \quad\text{for every }i,j.      \tag{B-RD64.1}
\]
Define
\[
 q_0=1,\qquad q_{j+1}=\min\{a+q_j,\ p q_j+p-1\}.
                                                               \tag{B-RD64.2}
\]
If $f:F\to G$ has no nonzero finite flat subgroup in its kernel,
has finite special-fibre kernel, and $q_u\ge a/(p-1)$, then
\[
                         F^u\ker(f_k)=0.              \tag{B-RD64.3}
\]
Here $F^u$ is iterated relative Frobenius. The assumptions and
coefficient bound concern this actual target $G$, not a source
or a group merely isogenous to $G$.
\end{lemma}
\begin{proof}
Lift an eigenbasis of the special cotangent space to parameters,
and average each parameter using
$N^{-1}\sum_{j=0}^{N-1}\zeta^{-jc_i}\tau^j$.
The reductions still form a basis, so the averaged parameters
are integral formal coordinates, with exact ring action
$\tau(X_i)=\zeta^{c_i}X_i$. This uses only that $N$ is a unit
and its Teichmuller roots lie in $K$. For a coefficient
$b_{j,\alpha}$ of $X^\alpha$ in the $j$th coordinate of $[p]_G$,
commutation with this action gives
\[
 \tau(b_{j,\alpha})=
       \zeta^{c_j-\sum_i\alpha_i c_i}b_{j,\alpha}.       \tag{B-RD64.4}
\]
The opposite point-action convention reverses both exponents
and gives the same ensuing divisibility. A nonzero eigen-coefficient
has valuation congruent to its exponent modulo $N$: divide it by
$\pi^{v_L(b_{j,\alpha})}$ and reduce the resulting unit, on which
$\tau$ acts trivially. For $|\alpha|=p$, (B-RD64.1) therefore gives
\[
                         v_L(b_{j,\alpha})\equiv0\pmod{p-1}.
\]
The product-height-two special-fibre condition makes every
coefficient of total degree less than $p^2$ vanish modulo $\pi$.
In particular every nonzero degree-$p$ coefficient has valuation
at least $p-1$.

There is a second coefficient constraint, independent of the
chosen eigenbasis. If $U(X)dX$ is the matrix of normalized
invariant differentials, then
\[
                         U([p]X)J_{[p]}(X)=pU(X).
\]
Thus a monomial with any exponent not divisible by $p$ has
coefficient divisible by $p$. The linear term is $pX$.
Apart from the degree-$p$ monomials, every nonlinear monomial
of total degree less than $2p$ has an exponent not divisible
by $p$. The remaining terms of degree $2p$ through $p^2-1$
have coefficients in $\pi R$, and all later coefficients are
integral. Consequently an augmentation vector of coefficient
valuation at least $q\ge1$ is carried by $[p]_G$ into one of
coefficient valuation at least
\[
 \min\{a+q,\ p-1+pq,\ 1+2pq,\ p^2q\}
             =\min\{a+q,\ p-1+pq\}.                   \tag{B-RD64.5}
\]
The last equality follows from $q\ge1$ and $p\ge5$.
The special order condition is preserved by the averaging
coordinate change: substitutions with invertible linear part
preserve the ideal of total degree at least $p^2$.
Iteration proves the valuations in (B-RD64.2) for
$[p]^j_G(\pi h)$, $h\in tR[[t]]^d$.

For clarity, we give the exact passage to the kernel estimate.
The coefficient $l_\alpha$ of a normalized logarithm, with
$n=|\alpha|>0$, satisfies
\[
 v_L(l_\alpha)\ge-a\lfloor\log_p n\rfloor,
 \qquad \lfloor\log_p n\rfloor\le\frac{n-1}{p-1}.
\]
This follows by integrating the invariant differentials, exactly
as in Lemma~\ref{lem:B-RD34-log}. If $q_u\ge a/(p-1)$, formal
substitution, coefficient by coefficient, therefore gives
\[
               p^u\log_G(\pi h)\in\pi^{q_u}R[[t]]^d.  \tag{B-RD64.6}
\]
Equality at the bound is permitted: each coefficient of an
augmentation-series substitution involves only finitely many terms.

Set $J=\pi^{\lceil a/(p-1)\rceil}R$ and
$\mathcal R=R[[\Delta]][1/p]$. The original Cartier identities are
\[
 \mathcal C_G=\Omega\pi D_{G_\pi},\qquad
 J[[\Delta]]^d\subset D_G,\qquad
 E(G_k)=D_G/(D_G\cap\mathcal R^d).
\]
They are Proposition~10(4), Proposition~9, and Lemma~3 of
\cite[pp.~674--676]{Bondarko2006Generic}; the bounded-projection
identity is Proposition~11(2). Integral Witt-coefficient action
of $\Omega$ preserves every coefficient-valuation ideal, so
(B-RD64.6) implies $p^u\mathcal C_G\subset J\{\{\Delta\}\}^d$.
For $z\in D_G$ and an integral matrix $A$ with $Az\in\mathcal R^b$,
the bounded-projection identity and
$\mathbf fz=p(z-z(0))\Delta^{-1}$ imply
\[
                         \mathbf f^uAz\in J[[\Delta]]^b.
\]
The discarded constant terms lie in $pR\subset J$, exactly as
in (B-RD34.8).

Finally choose an integral complement to
$df(0)(L^{\dim F})\cap R^d$ and its projection $\operatorname{pr}$.
Absence of a nonzero finite flat kernel gives
$df(0)D_F=D_G\cap df(0)(L^{\dim F})[[\Delta]]$.
For a class in the saturation of $E(F_k)$ in $E(G_k)$,
represented by $z\in D_G$, the projected series is bounded.
Thus $\mathbf f^u\operatorname{pr}z\in J[[\Delta]]^d\subset D_G$
is bounded and reduces to zero, while its difference from
$\mathbf f^uz$ belongs to $df(0)D_F$.
The saturation quotient is killed by $\mathbf f^u$.
Proposition~16 with its full proof, p.~681, and Theorem~2(8),
equation~(6), p.~668 of the cited source identify this quotient
with the Dieudonne module of $\ker f_k$, with $\mathbf f$
equal to relative Frobenius. This proves (B-RD64.3).
\end{proof}

\begin{proposition}[Common determinants control the quotient weights]
\label{prop:B-RD64-formal}
\status{Candidate; complete graph, pairing, and weight proof written.}
Let $p\ge5$ and let $K$ be complete discretely valued of mixed
characteristic $(0,p)$ with algebraically closed residue field
and absolute index $e$. Let $M,N$ be height-two groups with
potentially supersingular good reduction and common determinant
$\chi_p\psi$, where $\psi$ is a finite $\mathbf Z_p^\times$-valued character.
Put $a=e(p^2-1)$ and define $q_j$ by (B-RD64.2).
If $n\ge2$ and $q_{n-1}\ge a/(p-1)$, then an isomorphism
$M[p^n]\simeq N[p^n]$ determines their entire finite special
inertia characters.
\end{proposition}
\begin{proof}
Assume instead that the special characters differ. By
Lemma~\ref{lem:B-RD49-inertia}, both groups have good formal
models over a common tame field of degree dividing $N_0=p^2-1$.
Enlarge it to the unique tame extension $L/K$ of degree $N_0$;
this extension is not required to act faithfully on either
special group. Choose $\pi_L$ with $\pi_L^{N_0}=\pi_K$ and
$\tau\pi_L=\zeta\pi_L$.
Let $M_0,N_0'$ denote the two special formal groups.
Their special eigenvalue multisets are Frobenius orbits
$\{\lambda,\lambda^p\}$ and $\{\mu,\mu^p\}$ in
$\mu_{p^2-1}$, allowing repeated roots. Distinct orbits are
disjoint. The characteristic polynomials over $\mathbf Z_p$
therefore have relatively prime reductions, so an integral
polynomial in $\tau$ projects $M_0\times N_0'$ onto either factor.

Let $H$ be the finite-flat closure of the given generic graph
inside $\mathcal M[p^n]\times\mathcal N[p^n]$.
Over $L$ both Tate determinants are cyclotomic, because the finite
character is killed on the good field. The alternating generic pairings
and their inverse duality maps extend to perfect pairings on
the full good BT models by
\cite[Theorem~4 and Corollaries~1--2, p.~181]{Tate1967PDivisible}.
Scale one pairing by the unit multiplier of the torsion map.
Then $H$ is its own orthogonal: it is isotropic by flatness,
its orthogonal is finite flat of the same rank, and the
equal-rank closed immersion is an isomorphism. This identity
commutes with special-fibre base change by exact finite-flat
Cartier duality, as in the full proof of Proposition~\ref{prop:B-RD33}.

The integral tame projector preserves $H_k$. Hence
$H_k=J_M\times J_N$. The unique subgroup of rank $p^r$ in a
one-dimensional formal group over $k$ is $\ker F^r$:
its ideal is $(T^{p^r})$ in $k[[T]]$. Self-orthogonality in
the $p^n$-torsion, of rank $p^{2n}$ in each factor, forces
both subgroup exponents to equal $n$. Therefore
\[
                     H_k=\ker F_{M_0}^n\times
                                  \ker F_{N_0'}^n.      \tag{B-RD64.7}
\]

Set $Q=(\mathcal M\times\mathcal N)/H$. Quotient formation
commutes with base change, so its special fibre is the actual
product $M_0^{(p^n)}\times (N_0')^{(p^n)}$, with its induced
descent action. This is an isomorphism of special groups with
descent, not just an isogeny. Put $G=Q\times\mathcal M$.
The map
\[
 \mathcal M\times\mathcal N\longrightarrow G,
                \qquad z\longmapsto(qz,\operatorname{pr}_{\mathcal M}z)
                                                               \tag{B-RD64.8}
\]
has no nonzero finite flat subgroup in its kernel, since its
generic graph meets $0\times N$ trivially. Its special kernel
is exactly $0\times\ker F_{N_0'}^n$.

We now verify the crucial hypothesis on this actual target.
Write the special cotangent eigenvalues of $M_0,N_0'$ as
$\bar\zeta^{c_M},\bar\zeta^{c_N}$. Their reduced norms are the
same finite determinant value $\psi(\tau)$, by the covariant
dictionary and determinant calculation in
Proposition~\ref{prop:B-RD4-dictionary} and
Lemma~\ref{lem:B-RD49-inertia}. Thus
$c_M\equiv c_N\pmod{p-1}$. A simultaneous inversion convention
would give the same congruence.
Relative Frobenius of order $n$ raises cotangent characters
to their $p^n$th powers. Consequently the three eigen-exponents
of $G_k$ are
\[
                          p^n c_M,\quad p^n c_N,\quad c_M
                          \pmod{p^2-1}.
\]
They are all congruent modulo $p-1$. Hence
Lemma~\ref{lem:B-RD64-weighted-kernel} applies to (B-RD64.8)
with $u=n-1$. It would kill $\ker F_{N_0'}^n$ by $F^{n-1}$.
Its image instead has rank $p$, a contradiction. This excludes
different special characters under the stated bound. No
thickness of $Q$ was inferred from thickness of its sources;
the degree-$p$ coefficient bound was proved from $Q$'s own
semilinear descent action.
\end{proof}

\begin{theorem}[The next cyclotomic ramification has least index two]
\label{thm:B-RD64}
\status{Candidate; complete reduction and exact-index proof written.}
For every prime $p\ge5$, let $K$ range over complete discretely
valued mixed-characteristic $(0,p)$ fields in the full B.4
monodromy setting, with height two, $\mu_p\subset K$, perfect
residue field and
$e_K=2(p-1)$,
\[
                         \nu_{\min}(p,2(p-1),2)=2.     \tag{B-RD64.9}
\]
In particular the actual $p^3$-torsion always determines the
requested character, while actual $p^2$-torsion need not.
\end{theorem}
\begin{proof}
The lower bound is the already proved Theorem~\ref{thm:B-RD62}
with $n=2$, since $p<2(p-1)<p^2$.
For sufficiency suppose two allowed groups have isomorphic
$p^3$-torsion, and pass to the completed maximal unramified
base. Their generic determinants have the form
$\chi_p^d\epsilon$ with $d\in\{0,1,2\}$ and finite
$\mu_{p-1}$-valued $\epsilon$. The cyclotomic image modulo
$p^3$ has order $p^2$, because $\mu_p\subset K$ and
$p\nmid e_K$. Comparing determinants therefore forces the
same $d$ (their difference has absolute value at most two),
and then the same finite determinant character.

For $d=0$ the special character is the Brauer lift of the
semisimplified generic mod-$p$ inertia representation, with
additivity across the monodromy filtration. Cartier duality
reduces $d=2$ to this case. For $d=1$, each group is either
wholly potentially supersingular or is filtered by two
potentially good height-one groups: the connected--etale
filtration of a potentially ordinary graded piece descends,
as in the full proof of Theorem~\ref{thm:B-RD49}.
Two filtered groups have their finite special characters
determined by the semisimplified generic mod-$p$ representation,
since $\chi_p\equiv1\pmod p$ on inertia.

A mixed comparison is impossible by
Proposition~\ref{prop:B-RD63-mixed} at $n=3$, because
$p\nmid2(p-1)$ and $2(p-1)<p^3$.
It remains to compare two whole potentially supersingular
groups with common determinant. For their enlarged common
tame good field take
\[
 a=2(p-1)(p^2-1),\qquad q_0=1,\quad q_1=2p-1,
 \quad q_2=2p^2-1>2(p^2-1)=a/(p-1).
\]
These are exactly the first two values of (B-RD64.2).
Proposition~\ref{prop:B-RD64-formal} with $n=3$ excludes the
last possible character difference. Thus manuscript index
$\nu=2$ always suffices, and the lower example proves equality.
\end{proof}

\begin{remark}[Scope of the new target argument]
\label{rem:B-RD64-scope}
\status{Candidate scope statement.}
The coefficient lemma proves additional divisibility on the
actual quotient target by its tame action. It does not assume
that a ramified quotient preserves the Hasse-coefficient
valuation of a source. The complete exact-index assertion is
only for $p\ge5$, $e=2(p-1)$ and height two. The formal-pair
proposition has its own explicit recurrence bound and common
determinant hypothesis; it is not an unrestricted least-index
formula for other heights or ramification indices.

The exact supplement
\path{code/reduction_descent/verify_tame_quotient_weights.py}
checks $78770$ distinct equal-norm orbit pairs, $236310$ actual
quotient-weight triples, $2126790$ degree-$p$ coefficient bounds,
and $1740891$ low-degree monomials for nine primes. It also checks
$1800$ recurrence identities, $45000$ logarithm inequalities and
$1485$ orthogonality cases. The author run passed in
$1.6924098020026577$ seconds at 20:13:57 UTC. Its output retains
the exact scope; these finite checks prove no geometric existence
or universal formal-group theorem.
\end{remark}
